% !TEX program = xelatex
\documentclass[12pt,a4paper]{article}
\usepackage[fontset=windows]{ctex}
\usepackage{geometry}
\geometry{left=2.5cm,right=2.5cm,top=2.5cm,bottom=2.5cm}
\usepackage{amsmath}
\usepackage{amsfonts}
\usepackage{graphicx}
\usepackage{booktabs}
\usepackage{caption}
\usepackage{subcaption}
\usepackage{siunitx}
\usepackage{pgfplots}
\usepackage{pgfplotstable}
\usepackage{xurl} % 允許長檔名或URL自動換行避免超出邊界

\pgfplotsset{compat=1.18}
\usepgfplotslibrary{groupplots}
\sisetup{round-mode=places, round-precision=3}

% 定義 CSV 檔案路徑（必須與 .tex 在同一目錄�?
\pgfplotstableread[col sep=comma]{2d_ssss_vibration_q4_5x5_mindlin_exact_comparison.csv}{\dataFive}
\pgfplotstableread[col sep=comma]{2d_ssss_vibration_q4_17x17_mindlin_exact_comparison.csv}{\dataSeventeen}
\pgfplotstableread[col sep=comma]{2d_ssss_vibration_q4_5x5_17x17_selected_modes.csv}{\dataSelected}

% 頁面設定
\raggedbottom
\allowdisplaybreaks

\begin{document}

\title{簡支方板（SSSS）自由振動有限元與Mindlin解析解對比報告}
\author{基於 Q4 單元5×5 17×17 網格}
\date{}
\maketitle

\section{目的}

本報告整理自由振動結果
使用無因次頻率 \(\Omega\) 作為比較指標，定義如下：

\[
\Omega = \omega a^2 \sqrt{\frac{\rho h}{D}}, \qquad 
D = \frac{E h^3}{12(1-\nu^2)}
\]

比較 5×5 和 17×17 網格的有限元數值解SSSS Mindlin/Reissner 厚板解析解
使用 \(\Omega\) 可以消除尺寸、密度與彎曲剛度的量綱影響，比直接比較 Hz 更適合做 benchmark 指標?

\section{無因次頻率折線圖}

\ref{fig:omega-5x5} \ref{fig:omega-17x17} 分別比較 5×5 和 17×17 網格的有限元無因次頻率 \(\Omega_{\mathrm{FEM}}\) 與解析解 \(\Omega_{\mathrm{exact}}\)?
由於 5×5 網格在高階模態出現非常大的跳躍，兩圖均採用對數刻度?

\begin{figure}[htbp]
\centering
\begin{tikzpicture}
\begin{axis}[
    xlabel={Mode rank},
    ylabel={$\Omega$},
    ymode=log,
    log basis y=10,
    grid=major,
    legend pos=north west,
    width=0.9\textwidth,
    height=0.45\textwidth,
    title={5×5 網格FEM 頻率與解析解對比},
    mark size=1.5pt
]
\addplot[smooth,mark=o,color=blue] table[x=mode_rank, y=Omega_exact] {\dataFive};
\addlegendentry{解析解$\Omega_{\mathrm{exact}}$}

\addplot[smooth,mark=square,color=red] table[x=mode_rank, y=Omega_num] {\dataFive};
\addlegendentry{5×5 網格 $\Omega_{\mathrm{FEM}}$}
\end{axis}
\end{tikzpicture}
\caption{5×5 網格的無因次頻率對比}
\label{fig:omega-5x5}
\end{figure}

\begin{figure}[htbp]
\centering
\begin{tikzpicture}
\begin{axis}[
    xlabel={Mode rank},
    ylabel={$\Omega$},
    ymode=log,
    log basis y=10,
    grid=major,
    legend pos=south east,
    width=0.9\textwidth,
    height=0.45\textwidth,
    title={17×17 網格FEM 頻率與解析解對比},
    mark size=1.5pt
]
\addplot[smooth,mark=o,color=blue] table[x=mode_rank, y=Omega_exact] {\dataSeventeen};
\addlegendentry{解析解$\Omega_{\mathrm{exact}}$}

\addplot[smooth,mark=square,color=red] table[x=mode_rank, y=Omega_num] {\dataSeventeen};
\addlegendentry{17×17 網格 $\Omega_{\mathrm{FEM}}$}
\end{axis}
\end{tikzpicture}
\caption{17×17 網格的無因次頻率對比}
\label{fig:omega-17x17}
\end{figure}


\section{無因次頻率相對誤差折線圖}

定義相對誤差百分比為

\[
\epsilon_{\Omega} = \frac{\Omega_{\mathrm{FEM}} - \Omega_{\mathrm{exact}}}{\Omega_{\mathrm{exact}}} \times 100\%.
\]

\ref{fig:error-5x5} \ref{fig:error-17x17} 分別展示兩種網格的 \(|\epsilon_{\Omega}|\)（對數縱軸），以便觀察誤差的量級變化
5×5 網格在低階已有約 70\% 的誤差，高階誤差急劇增大；17×17 網格在整個計算範圍內誤差均低於 20\%

\begin{figure}[htbp]
\centering
\begin{tikzpicture}
\begin{axis}[
    xlabel={Mode rank},
    ylabel={$|\epsilon_{\Omega}|$ (\%)},
    ymode=log,
    log basis y=10,
    grid=major,
    legend pos=north west,
    width=0.9\textwidth,
    height=0.45\textwidth,
    title={5×5 網格的相對誤差絕對值},
    mark size=1.5pt
]
\addplot[smooth,color=red] table[x=mode_rank, y=abs_rel_error_percent] {\dataFive};
\addlegendentry{5×5 網格}
\end{axis}
\end{tikzpicture}
\caption{5×5 網格的相對誤差絕對值}
\label{fig:error-5x5}
\end{figure}

\begin{figure}[htbp]
\centering
\begin{tikzpicture}
\begin{axis}[
    xlabel={Mode rank},
    ylabel={$|\epsilon_{\Omega}|$ (\%)},
    ymode=log,
    log basis y=10,
    grid=major,
    legend pos=north east,
    width=0.9\textwidth,
    height=0.45\textwidth,
    title={17×17 網格的相對誤差絕對值},
    mark size=1.5pt
]
\addplot[smooth,color=green] table[x=mode_rank, y=abs_rel_error_percent] {\dataSeventeen};
\addlegendentry{17×17 網格}
\end{axis}
\end{tikzpicture}
\caption{17×17 網格的相對誤差絕對值}
\label{fig:error-17x17}
\end{figure}

\section{指定模態數據表}

\ref{tab:selected-5x5} \ref{tab:selected-17x17} 分別摘錄 5×5 和 17×17 網格中部分模態的無因次頻率與誤差。
表中的 \(m\) 與 \(n\) 分別代表 Mindlin 板解析解中 \(x\) 方向與 \(y\) 方向的半波數（Half-wave numbers），即在該方向上有幾個波峰。
其中 1 表示殘差通過，0 表示未通過，殘差容差設定為 \(\text{tol} = 1 \times 10^{-6}\)。

\vspace{1cm}

\begin{table}[htbp]
\centering
\caption{表A: 5×5 網格\(\Omega\) 對比表}
\label{tab:selected-5x5}
% 使用 resizebox 避免表格超出頁面邊界
\resizebox{\textwidth}{!}{
\pgfplotstabletypeset[
    columns={mode_rank, m, n, Omega_exact, Omega_num_5x5, rel_error_percent_5x5, residual_ok_5x5},
    columns/mode_rank/.style={column name=Mode},
    columns/m/.style={column name=$m$},
    columns/n/.style={column name=$n$},
    columns/Omega_exact/.style={column name=$\Omega_{\text{exact}}$, sci, sci zerofill, precision=3},
    columns/Omega_num_5x5/.style={column name=$\Omega_{\text{FEM}}$, sci, sci zerofill, precision=3},
    columns/rel_error_percent_5x5/.style={column name=誤差 (\%), fixed, fixed zerofill, precision=1},
    columns/residual_ok_5x5/.style={column name=殘差},
    every head row/.style={before row=\toprule, after row=\midrule},
    every last row/.style={after row=\bottomrule},
    string type
]{\dataSelected}
}
\end{table}

\begin{table}[htbp]
\centering
\caption{表B: 17×17 網格\(\Omega\) 對比表}
\label{tab:selected-17x17}
\resizebox{\textwidth}{!}{
\pgfplotstabletypeset[
    columns={mode_rank, m, n, Omega_exact, Omega_num_17x17, rel_error_percent_17x17, residual_ok_17x17},
    columns/mode_rank/.style={column name=Mode},
    columns/m/.style={column name=$m$},
    columns/n/.style={column name=$n$},
    columns/Omega_exact/.style={column name=$\Omega_{\text{exact}}$, sci, sci zerofill, precision=3},
    columns/Omega_num_17x17/.style={column name=$\Omega_{\text{FEM}}$, sci, sci zerofill, precision=3},
    columns/rel_error_percent_17x17/.style={column name=誤差 (\%), fixed, fixed zerofill, precision=1},
    columns/residual_ok_17x17/.style={column name=殘差},
    every head row/.style={before row=\toprule, after row=\midrule},
    every last row/.style={after row=\bottomrule},
    string type
]{\dataSelected}
}
\end{table}

\section{簡要解讀}

\begin{itemize}
    \item \textbf{5×5 網格}：前 50 個模態的 \(\Omega_{\mathrm{FEM}}\) 顯著高於 Mindlin 解析解（相對誤差 70\% 00\%），表明離散模型整體偏硬。第 60 階以後有限元頻率急劇增大數個量級，對應非物理的剪切鎖定或數值噪聲，這些模態不宜作為主要彎曲模態解讀
    \item \textbf{17×17 網格}：在整個計算的 500+ 階模態中，\(\Omega_{\mathrm{FEM}}\) 與解析解吻合良好，相對誤差始終低 m20\%，多數模態誤差小 10\%。殘差檢查（欄位 \texttt{residual\_ok}）顯示17×17 網格除前幾個模態外均滿足容差，5×5 網格在低階未通過（殘差容差設置較寬鬆時仍可能標記 0，實際需根據程序設定判斷）
    \item \textbf{建議}：對於定量基準測試，推薦使用至少 17×17  Q4 網格5×5 網格僅適合趨勢預分析，不能用於準確頻率計算
\end{itemize}

\end{document}

